% !TeX root = ../main.tex
% Section 7: Discussion and Policy Implications
% Batch #4.7 — Post M1+M2 axis, notation hygiene, honest Tab 6 positioning

\section{Discussion and Policy Implications}\label{sec:discussion}

The multi-event lead structure of the quantum coherence signal $\mathcal{C}(\rho)$ over
conventional Pearson correlation (Section~\ref{M-sec:coherence-2023})---ranging
from contemporaneous for sudden failures to several weeks for gradually developing
ones---and the
15-trading-day lead of the von Neumann entropy $S(\rho)$ over the SVB failure
event (Section~\ref{M-sec:entropy-2023}), carry theoretical and policy implications
examined below.

\subsection{The Economic Planck Constant as a Microstructural Liquidity Measure}
\label{sec:hbar-liquidity}

$\hbar_{\mathrm{econ}}$ offers a structural liquidity measure grounded in market
microstructure theory. Unlike bid--ask spreads or volume-based metrics,
$\hbar_{\mathrm{econ}}$ admits a direct interpretation as $\lambda \cdot q_0$,
linking it to Kyle's price impact parameter. In our calibration
(Section~\ref{M-sec:data-methodology}), the full-sample mean of
$\hbar_{\mathrm{econ}}$ is 0.062, with time-series variation reflecting
regime-dependent liquidity conditions in the U.S.\ banking sector.

The measure complements conventional indicators in three respects: it derives from
a structural model, aggregates information across multiple observables (via the
Amihud illiquidity proxy applied to the ten-ticker equity/ETF panel), and connects
directly to option pricing via the effective volatility expansion
\eqref{M-eq:effective-vol}.

\subsection{Quantum Coherence as an Early Warning Indicator}
\label{sec:coherence-early-warning}

$\mathcal{C}(\rho)$ functions as an early warning indicator because the
non-commutative order-flow channel encodes correlations in $\rho$ before they
appear in price data. The commutation relation \eqref{M-eq:weyl-ccr} generates
institutional linkages through $\hat{P}$ that populate the density matrix ahead
of observable market stress. In our March~2023 case study
(Figure~\ref{M-fig:coherence-2023}), $\mathcal{C}(\rho_t)$ crosses its $+1\sigma$
band on March~9---the day \emph{before} SVB's collapse---while the classical
Pearson correlation between KRE and VIX crosses its $-1\sigma$ band on the same
day. For the First Republic Bank failure on May~1, $\mathcal{C}(\rho_t)$ crosses
its threshold on March~31, substantially before the Pearson threshold
(May~1). Averaged across the three
failure events, the coherence signal leads for gradually developing failures
and responds contemporaneously for sudden ones.

Real-time implementation is computationally feasible. The single-qubit
parametrization reduces the QST reconstruction \eqref{M-eq:mst} to a
three-dimensional convex problem solvable in milliseconds, making the framework
suitable for regulatory monitoring at daily or intraday frequencies.

\subsection{Policy Implications for Financial Stability}\label{sec:policy}

The findings carry implications for three areas of financial stability policy.

\paragraph{Macroprudential surveillance.}
The QST framework provides a theoretically grounded methodology in which
$\theta_t$ and $S(\rho)$ serve as systemic risk summary statistics, complementing
SRISK \citep{brownlees2016srisk} and CoVaR \citep{adrian2016covar}. In particular,
the von Neumann entropy's forward-looking property
(Proposition~\ref{M-prop:granger-entropy}) contributes information not contained in
lagged option-implied volatilities alone.

\paragraph{Market microstructure regulation.}
The economic Planck constant $\hbar_{\mathrm{econ}}$ provides a quantitative
measure of aggregate market depth. Elevated $\hbar_{\mathrm{econ}}$ during
crisis periods, driven by widening Amihud illiquidity across the banking sector,
signals impaired liquidity provision that may warrant regulatory attention.
Under the theoretical mapping $\hbar_{\mathrm{econ}} = \lambda \cdot q_0$,
sustained elevation corresponds to persistent widening of Kyle's price-impact
parameter across the panel.

\paragraph{Crisis management.}
The advance lead time of quantum coherence $\mathcal{C}(\rho)$ documented in
Section~\ref{M-sec:coherence-2023} provides a policy-relevant window for
pre-emptive intervention before systemic contagion registers in conventional
correlation-based indicators. For First Republic, the coherence lead
extends to several weeks in our sample (see the main paper §5.2 for the
precise count). For the SVB and Signature failures, both the coherence signal
and the classical Pearson benchmark cross their respective thresholds on March~9,
2023---providing the same same-day early warning as conventional indicators while
operating through a structurally distinct, non-commutative channel. In practical
terms, a coherence-monitoring dashboard crossing its $+1\sigma$ band would have
flagged the March~2023 crisis buildup on the closing prices of March~9, 2023.

\subsection{Limitations and Future Research}\label{sec:limitations}

The framework has several limitations worth noting.

\paragraph{Dimensionality of the Hilbert space.}
The single-qubit parametrization captures the distressed/non-distressed dichotomy
but cannot represent multi-region configurations or sector-level heterogeneity.
Extending to multi-qubit or continuous-variable representations is a natural
generalization, though it raises computational challenges from exponential
state-space growth and requires new identification strategies for higher-order
observables.

\paragraph{Data availability constraints.}
The FRED HY-OAS series (BAMLH0A0HYM2) is available in our sample only from
2023-05-22 onwards. For dates prior to this, our credit-stress axis $r_y$ is
set to zero, and the coherence measure $\mathcal{C}(\rho) = \sqrt{r_x^2 + r_y^2}$
reduces to $|r_x|$ during the pre-May-2023 portion of the crisis window. This
truncation attenuates the coherence signal in the earliest portion of the crisis;
extending the HY-OAS series backward via alternative data sources (Bloomberg,
Refinitiv) would strengthen the pre-May signal.

\paragraph{Multi-underlying calibration.}
The IV-surface calibration in Section~\ref{sec:iv-surface} uses a canonical
banking-crisis reference smile calibrated to KRE. Extending to JPM and SPY in
Table~\ref{tab:iv-calibration} incurs a projection error, since the reference
$(a, b)$ parameters were not re-optimized per underlying. A multi-underlying
joint calibration with underlying-specific smile parameters is deferred to
future work.

\paragraph{Hedging performance in single-path case study.}
The 20-day hedging case study (Table~\ref{tab:hedging-rmshe}) demonstrates a
marginal ($\approx 0.75\%$) improvement of the quantum optimal hedge over the
BSM baseline. This magnitude reflects the specific $(x_t, \tau_t)$ configuration
of the SVB stress path. Proposition~\ref{M-prop:quantum-optimal-hedge} provides the
closed-form basis for larger improvements when the option is deeper OTM and
$\tau$ shorter (i.e., when the correction term
$\hbar_{\mathrm{econ}} \Gamma_{\text{BS}} \gamma x/\tau$ dominates). A
cross-strike, cross-maturity hedging comparison over multiple stress episodes
would establish the empirical distribution of improvements more broadly.
The Lindblad master equation \eqref{M-eq:lindblad} is currently phenomenological.
Deriving the Lindblad operators from limit order book dynamics, market maker
behavior, and information asymmetry would strengthen the theoretical
foundations.

\paragraph{Extension to American and exotic derivatives.}
The framework currently treats European-style claims. Extending to American
options with optimal exercise and path-dependent derivatives is an important
direction for both theory and practice.

Each direction constitutes a research program in its own right. We offer the
present framework as a foundation for these extensions.
